\section{Readdressing control operands after a decoder change}
\label{sec:control-readdressing}

A control operand is a relative address in the cons-tree layout.  A
continue instruction of depth two before a rotation may need depth three
afterwards to reach the same next instruction occurrence.  Requiring
literal equality of the two operands would reject this valid
readdressing.  The appropriate finite certificate compares their actual
resolved destinations.

For a consumed-path register program, assign each instruction occurrence
the operational label
\[
 L(v)=\begin{cases}
  (\mathrm{inc},r),&v=\mathrm{CI}(r),\\
  (\mathrm{dec},r),&v=\mathrm{CD}(r),\\
  \mathrm{halt},&v=\mathrm{CH},\\
  \mathrm{control},&v=\mathrm{CB}(k)\text{ or }v=\mathrm{CC}(k).
 \end{cases}
\]
The last two primitive kinds preserve every register.  Their operands
still matter when computing their successor: the checker applies the
actual strip-and-increment rule for each break, the actual strip rule
for each continue, and the candidate's own decoder with its implicit
zero suffix.  The operands are omitted only from the subsequent label
comparison, after their destinations have been calculated.

\begin{theorem}[Operational control-graph certificate]
Let $P$ and $P'$ be finite acyclic decoders.  Suppose a bijection $f$
between their instruction occurrences preserves operational labels and
initial control, and maps each resolved successor to the corresponding
successor.  For decrements, zero and positive successors must match in
their respective order.  Then the two consumed-path register programs
have identical halting behavior from corresponding controls and equal
register vectors, for every natural register vector.  Finite halting
computations also have equal primitive-step counts and final registers.
\end{theorem}
\begin{proof}
Relate $(v,R)$ in $P$ to $(f(v),R)$ in $P'$.  An increment has the same
operand and therefore the same update.  A decrement reads the same
register, takes the same guard, and has the same update.  A control
instruction changes no register.  The successor premise gives related
output controls in each case.  Halt labels agree and are terminal in
both programs.  This is a deterministic one-instruction bisimulation;
induction on a finite run proves both directions and preserves its
length and final register vector.
\end{proof}

For \texttt{s10\_best\_004}, the independent verifier compares the
candidate to the already certified color-106 commuted seed using
depth-first occurrence number as $f$.  Both have 6,460 occurrences and
8,390 guarded successor edges.  Every operational label, every resolved
successor, and the initial occurrence match.  Exactly 115 literal
operands differ, all from $\mathrm{CC}(2)$ to $\mathrm{CC}(3)$; no
increment or decrement operand changes.  The candidate has 203 shared
pair nodes and maximum path depth 45.  Its source equivalence follows
from the theorem and composes with the earlier region-contraction proof.

The portable checker is
\texttt{code/audit/check\_reassociated\_graph.py}, invoked with the
explicit \texttt{--operational-controls} option.  Without that option it
continues to require literal equality of all instruction operands.  Its
JSON record lists every changed control occurrence and explicitly records
that literal operand equality is false while operational-label and
resolved-successor equality are true.

This source result does not establish the blank-tape target for the
emitted binary table.  The separate affine runtime checker validates all
8,391 guarded frame contracts at minimum header distance 64, but the
blank bootstrap of this table has not been calibrated.  It is retained
as a validated source and runtime experiment, not as a new blank-input
consistency machine.  A successful bootstrap search could close this gap
after exact full-tape replay and verification of the inherited logical
initialization hypotheses; no such audited prefix has been found here.
